\documentclass[10pt,aspectratio=169]{beamer}
% Preámbulo modular para presentaciones Beamer (Modelación Estadística)
% Diseño y paleta institucional del Tecnológico de Monterrey

\documentclass[aspectratio=169,xcolor={dvipsnames,table}]{beamer}

\usepackage[utf8]{inputenc}
\usepackage[spanish,mexico]{babel}
\usepackage{amsmath,amssymb,amsfonts,mathtools}
\usepackage{graphicx}
\usepackage{listings}
\usepackage{booktabs}
\usepackage{tikz}
\usetikzlibrary{matrix,arrows,decorations.pathmorphing,shapes.geometric,positioning}

% Paleta Institucional Tecnológico de Monterrey
\definecolor{TecAzul}{HTML}{0039A6}       % Azul TEC: color primario institucional
\definecolor{TecAzulOscuro}{HTML}{1A2E51} % Azul marino oscuro
\definecolor{TecRojo}{HTML}{EC2661}       % Rojo de acento (paleta miTec)
\definecolor{TecGris}{HTML}{646464}       % Gris neutro para comentarios y subtítulos
\definecolor{TecFondoBloque}{HTML}{F4F6F9}% Fondo suave para bloques

% Configuración de Tema Beamer
\usetheme{Boadilla}
\usecolortheme[named=TecAzulOscuro]{structure}
\setbeamercolor{alerted text}{fg=TecRojo}
\setbeamercolor{palette primary}{bg=TecAzulOscuro,fg=white}
\setbeamercolor{palette secondary}{bg=TecAzul,fg=white}
\setbeamercolor{palette tertiary}{bg=TecAzulOscuro!80!black,fg=white}
\setbeamercolor{title}{fg=TecAzulOscuro,bg=TecFondoBloque}
\setbeamercolor{frametitle}{fg=TecAzulOscuro,bg=TecFondoBloque}

% Bloques personalizados elegantes
\setbeamercolor{block title}{bg=TecAzulOscuro,fg=white}
\setbeamercolor{block body}{bg=TecFondoBloque,fg=black}
\setbeamercolor{block title alerted}{bg=TecRojo,fg=white}
\setbeamercolor{block body alerted}{bg=TecRojo!8,fg=black}
\setbeamercolor{block title example}{bg=TecAzul,fg=white}
\setbeamercolor{block body example}{bg=TecAzul!8,fg=black}

% Redondear bordes y añadir sombra sutil
\setbeamertemplate{blocks}[rounded][shadow=true]
\setbeamertemplate{navigation symbols}{}

% Pie de página limpio con número de diapositiva
\setbeamertemplate{footline}{
  \leavevmode%
  \hbox{%
  \begin{beamercolorbox}[wd=.7\paperwidth,ht=2.25ex,dp=1ex,left]{author in head/foot}%
    \hspace*{2ex}\insertshortauthor~~(\insertshortinstitute) \hfill \insertshorttitle
  \end{beamercolorbox}%
  \begin{beamercolorbox}[wd=.3\paperwidth,ht=2.25ex,dp=1ex,right]{date in head/foot}%
    \insertshortdate{}\hspace*{2em}
    \insertframenumber{} / \inserttotalframenumber\hspace*{2ex}
  \end{beamercolorbox}}%
  \vskip0pt%
}

% Configuración de Listings de Python para visualización pedagógica de código
\lstset{
    language=Python,
    basicstyle=\ttfamily\footnotesize,
    keywordstyle=\color{TecAzulOscuro}\bfseries,
    stringstyle=\color{TecRojo},
    commentstyle=\color{TecGris}\itshape,
    showstringspaces=false,
    breaklines=true,
    frame=lines,
    rulecolor=\color{TecAzulOscuro!30},
    backgroundcolor=\color{TecFondoBloque},
    aboveskip=0.5em,
    belowskip=0.5em,
    literate={á}{{\'a}}1 {é}{{\'e}}1 {í}{{\'i}}1 {ó}{{\'o}}1 {ú}{{\'u}}1
             {Á}{{\'A}}1 {É}{{\'E}}1 {Í}{{\'I}}1 {Ó}{{\'O}}1 {Ú}{{\'U}}1
             {ñ}{{\~n}}1 {Ñ}{{\~N}}1 {¿}{{?`}}1 {¡}{{!`}}1
}

% Metadatos institucionales por defecto
\title[Modelación Estadística]{Modelación Estadística para Ciencia de Datos e Ingeniería}
\author[J. Castillo Colmenares]{Juliho David Castillo Colmenares}
\institute[Tec de Monterrey]{Tecnológico de Monterrey \\ Escuela de Ingeniería y Ciencias}
\date{\today}
% Comandos y macros estadísticas compartidas para las presentaciones Beamer (Metropolis)

% Conjuntos numéricos básicos
\newcommand{\R}{\mathbb{R}}
\newcommand{\N}{\mathbb{N}}
\newcommand{\Q}{\mathbb{Q}}
\newcommand{\Z}{\mathbb{Z}}
\newcommand{\set}[1]{\left\{ #1 \right\}}
\newcommand{\abs}[1]{\left|#1\right|}
\newcommand{\norm}[1]{\left\| #1 \right\|}

% Comandos estadísticos y probabilísticos del curso
\newcommand{\Var}{\operatorname{Var}}
\newcommand{\cov}{\operatorname{Cov}}
\newcommand{\corr}{\rho}
\newcommand{\comb}[2]{\begin{pmatrix} #1 \\ #2 \end{pmatrix}}
\newcommand{\s}{\sigma}
\newcommand{\particion}{\mathfrak{P}}
\newcommand{\card}[1]{\#\left( #1 \right)}
\newcommand{\seq}[3]{\set{#1}_{#2}^{#3}}
\newcommand{\E}{\mathbb{E}}
\newcommand{\Prob}{\mathbb{P}}

% Entornos pedagógicos y cajas en español académico natural
\newenvironment{discusion}{%
  \begin{alertblock}{Pregunta de Discusión}%
}{%
  \end{alertblock}%
}

\newenvironment{reflexion}{%
  \begin{alertblock}{Reflexión para el Aula}%
}{%
  \end{alertblock}%
}

\newenvironment{paradoja}{%
  \begin{alertblock}{Paradoja de Exploración}%
}{%
  \end{alertblock}%
}

\newenvironment{motivacion}{%
  \begin{exampleblock}{Motivación: Ciencia de Datos en la Práctica}%
}{%
  \end{exampleblock}%
}

\newenvironment{retoclase}{%
  \begin{block}{Reto Rápido en Equipo}%
}{%
  \end{block}%
}

\title[Efectos de modelo fijo]{Sección 08.03: Análisis de efectos de modelo fijo}\subtitle{Cuadrados medios y estadístico $F$}
\author[J. Castillo Colmenares]{Juliho Castillo Colmenares}\institute{Tecnológico de Monterrey}\date{}
\begin{document}
\begin{frame}[plain]\titlepage\end{frame}
\begin{frame}{Hoja de ruta}\small\begin{enumerate}\item Cuadrados medios.\item Hipótesis y estadístico $F$.\item Tabla ANOVA e interpretación.\end{enumerate}\end{frame}
\begin{frame}{Fuente y alcance}\small La sección analiza cómo los efectos de tratamiento cambian los cuadrados medios de un ANOVA de un factor.\begin{block}{Pregunta guía}¿Por qué el cociente entre cuadrados medios detecta efectos?\end{block}\end{frame}
\begin{frame}{Cuadrados medios}\small\[\text{CMTR}=\frac{\text{SCTR}}{k-1},\qquad \text{CME}=\frac{\text{SCE}}{N-k}.\]Cada suma de cuadrados se divide por sus grados de libertad.\end{frame}
\begin{frame}{Esperanzas bajo el modelo}\small\[E(\text{CME})=\sigma^2,\qquad E(\text{CMTR})=\sigma^2+\frac{\sum_i n_i\tau_i^2}{k-1}.\]\end{frame}
\begin{frame}{Hipótesis omnibus}\small\[H_0:\tau_1=\cdots=\tau_k=0\quad\Longleftrightarrow\quad\mu_1=\cdots=\mu_k.\]\[H_a:\text{al menos un efecto no es cero}.\]\end{frame}
\begin{frame}{Lógica del cociente}\small Bajo $H_0$, CMTR y CME estiman la misma varianza. Bajo $H_a$, CMTR incorpora un término positivo asociado a los tratamientos.\end{frame}
\begin{frame}{Estadístico $F$}\small\[F=\frac{\text{CMTR}}{\text{CME}}=\frac{\text{SCTR}/(k-1)}{\text{SCE}/(N-k)}.\]\end{frame}
\begin{frame}{Distribución nula}\small Bajo los supuestos del modelo y $H_0$,\[F\sim F_{k-1,N-k}.\]La distribución es central cuando todos los efectos son nulos.\end{frame}
\begin{frame}{Regla de decisión}\small Rechazar $H_0$ al nivel $\alpha$ si\[F>F_{\alpha,k-1,N-k}\]o, equivalentemente, si el valor $p$ es menor que $\alpha$.\end{frame}
\begin{frame}{Tabla ANOVA}\small\begin{itemize}\item Tratamientos: SCTR, $k-1$, CMTR.\item Error: SCE, $N-k$, CME.\item Total: SCT, $N-1$.\end{itemize}\end{frame}
\begin{frame}{Procedimiento I}\small\begin{enumerate}\item Revisar aleatorización, independencia y normalidad de errores.\item Calcular medias y suma de cuadrados.\item Obtener CMTR y CME.\end{enumerate}\end{frame}
\begin{frame}{Procedimiento II}\small\begin{enumerate}\setcounter{enumi}{3}\item Formar $F$.\item Comparar con el cuantil o calcular $p$.\item Reportar la conclusión global.\end{enumerate}\end{frame}
\begin{frame}{Aplicación: tratamientos}\small Un experimento compara $k$ tratamientos con $n_i$ réplicas. Las diferencias entre medias alimentan SCTR; la variación interna alimenta SCE.\end{frame}
\begin{frame}{Aplicación: medias}\small La media general $\bar Y_{..}$ y las medias de tratamiento $\bar Y_{i.}$ separan variación entre tratamientos de variación dentro de tratamientos.\end{frame}
\begin{frame}{Aplicación: cociente}\small Si CMTR es mucho mayor que CME, el cociente $F$ crece y la evidencia contra la igualdad de medias aumenta.\end{frame}
\begin{frame}{Aplicación: decisión}\small Un rechazo global afirma que al menos una media difiere; todavía no identifica qué pares son diferentes.\end{frame}
\begin{frame}{Supuestos}\small La validez depende de errores independientes, varianza común y normalidad aproximada. La aleatorización sustenta la independencia.\end{frame}
\begin{frame}{Interpretación}\small El resultado es una afirmación global sobre tratamientos. La magnitud práctica debe acompañarse con medias, diferencias e intervalos.\end{frame}
\begin{frame}{Comunicación}\small Reporte SCTR, SCE, grados de libertad, cuadrados medios, $F$, valor $p$ y la conclusión contextual.\end{frame}
\begin{frame}{Síntesis}\small\begin{itemize}\item CME estima el ruido.\item CMTR combina ruido y efectos.\item $F$ compara ambos.\item El contraste es global.\end{itemize}\end{frame}
\begin{frame}{Conexión con el capítulo}\small Tras un $F$ significativo se requieren comparaciones múltiples; la siguiente sección presenta LSD, Tukey y Bonferroni.\end{frame}
\end{document}
